%% ma: L12-B15-0022
%% lop: 12
%% chuong: 5
%% bai: 15
%% dang: 12.15.18
%% loai: TLN
%% muc: VD
%% dap_an: "2,83"
%% nguon: "(NBV)-10 ĐỀ KIỂM TRA HK2-TOÁN 12-DUNG CHUNG 3 SÁCH.pdf (H03, Đề số 3), Phần 3 Câu 2"
%% kien_thuc: "Bài 15 (phương trình tham số của đường thẳng), Bài 8 (độ dài vectơ), hàm số bậc hai (lớp 10)"
%% trang_thai: chinh-thuc
%% ghi_chu: "Đáp án gốc 2,83 khớp. Kiểm bằng công thức khoảng cách giữa hai đường thẳng chéo nhau: 8/(2√2) = 2√2. Bản kho giữ cách giải bằng tổng bình phương (không dùng công thức khoảng cách giữa hai đường chéo nhau, ngoài SGK KNTT). Gốc mất dấu trừ ở phương trình d2 (đọc từ ảnh trang: (y-5)/(-2) = (z+3)/2). Đã duyệt 21/10/2026 theo yêu cầu của giáo viên: giải lại độc lập, kiểm bằng sympy, soát trùng, biên dịch OK."
\begin{ex}
Trong không gian $Oxyz$, cho hai đường thẳng chéo nhau
\[d_1\colon\begin{cases}x=1+t\\y=2\\z=4\end{cases}\qquad\text{và}\qquad d_2\colon\dfrac x1=\dfrac{y-5}{-2}=\dfrac{z+3}{2}.\]
Hai điểm $M$, $N$ lần lượt di chuyển trên $d_1$, $d_2$. Độ dài của đoạn thẳng $MN$ nhỏ nhất bằng bao nhiêu? (kết quả viết dưới dạng số thập phân và làm tròn đến hàng phần trăm).
\shortans{2,83}
\loigiai{$M(1+t;2;4)\in d_1$, $N(s;5-2s;2s-3)\in d_2$ nên $\overrightarrow{MN}=(s-t-1;3-2s;2s-7)$ và
\[MN^2=(s-t-1)^2+(3-2s)^2+(2s-7)^2=(s-t-1)^2+8\left(s-\dfrac52\right)^2+8\ge8.\]
Dấu bằng đạt khi $s=\dfrac52$, $t=\dfrac32$. Vậy $MN_{\min}=2\sqrt2\approx2{,}83$.}
\end{ex}
